\chapter{Cálculos de estabilidad}
\label{ap:estabilidad}

En este apéndice se presentan los cálculos detallados de estabilidad lineal para los 11 escenarios definidos en \texttt{scenarios.json}. Para cada escenario se calcula el espectro lineal $\lambda(k)=\mathrm{eig}(J-Dk^2)$, donde $J$ es el Jacobiano de reacción evaluado en el estado estacionario $(c^*,s^*,i^*)$ y $D=\mathrm{diag}(D_c,D_s,D_i)$ es la matriz de difusión.

\section{Allee débil}

Se linealiza el sistema espacial
\begin{align}
    \frac{dc}{dt} &= D_c \nabla^2 c + r_c c \left( \frac{c}{a} - 1 \right)\left(1 - \frac{c}{k_c} \right) - \alpha c s^2 - \beta c i^2 - \mu ( \gamma s^2 + \eta i^2), \label{eqn:cancer_dynamic_spatial_app} \\
    \frac{ds}{dt} &= D_s \nabla^2 s + r_s s \left(1 - \frac{s}{k_s} \right) - \gamma s c^2 + \delta s i^2 - \frac{\mu}{2} \alpha s c^2, \label{eqn:sano_dynamic_spatial_app} \\
    \frac{di}{dt} &= D_i \nabla^2 i + r_i i \left(1 - \frac{i}{k_i} \right) - \eta c^2 i + \delta s^2 i - \frac{\mu}{2} \beta i c^2. \label{eqn:inmune_dynamic_spatial_app}
\end{align}
Para un estado estacionario $(c_0, s_0, i_0)$, perturbaciones $\tilde{c},\tilde{s},\tilde{i}$ satisfacen
\[
\frac{d}{dt}
\begin{pmatrix}
 \tilde{c} \\
 \tilde{s} \\
 \tilde{i}
\end{pmatrix}
= 
\mathbf{J}
\begin{pmatrix}
 \tilde{c} \\
 \tilde{s} \\
 \tilde{i}
\end{pmatrix},
\]
donde $\mathbf{J}$ es el Jacobiano evaluado en $(c_0,s_0,i_0)$. El signo de $\Re(\lambda_j(k))$ determina estabilidad; partes imaginarias indican oscilaciones (remisión/recurrencia).

\subsection{Casos sin control adaptativo}

Se presentan los espectros lineales para los 4 escenarios de Allee débil sin control adaptativo:

\begin{figure}[H]
    \centering
    \begin{minipage}{0.48\textwidth}
        \centering
        \includegraphics[width=\textwidth]{images/estabilidad_lineal_weak_mu0_uNo_bajo_umbral.png}
        \subcaption{$\mu=0$, bajo umbral}
    \end{minipage}\hfill
    \begin{minipage}{0.48\textwidth}
        \centering
        \includegraphics[width=\textwidth]{images/estabilidad_lineal_weak_mu0_uNo_sobre_umbral.png}
        \subcaption{$\mu=0$, sobre umbral}
    \end{minipage}\\[0.5em]
    \begin{minipage}{0.48\textwidth}
        \centering
        \includegraphics[width=\textwidth]{images/estabilidad_lineal_weak_mu1_uNo_bajo_umbral.png}
        \subcaption{$\mu=1$, bajo umbral}
    \end{minipage}\hfill
    \begin{minipage}{0.48\textwidth}
        \centering
        \includegraphics[width=\textwidth]{images/estabilidad_lineal_weak_mu1_uNo_sobre_umbral.png}
        \subcaption{$\mu=1$, sobre umbral}
    \end{minipage}
    \caption{Espectros lineales para Allee débil sin control adaptativo. En todos los casos $\Re(\lambda_1)>0$ cerca de $k=0$, confirmando inestabilidad. El régimen $\mu=1$ atenúa modos de alta frecuencia pero no estabiliza el sistema.}
    \label{fig:stability_weak_no_control}
\end{figure}

\subsection{Casos con control adaptativo}

Se presentan los espectros para los 2 escenarios de Allee débil con control adaptativo tipo Hill $u=u_{\max}\,H_{\mathrm{act}}(c;K_c,n_c)\,H_{\mathrm{inh}}(i;K_i,n_i)$:

\begin{figure}[H]
    \centering
    \begin{minipage}{0.48\textwidth}
        \centering
        \includegraphics[width=\textwidth]{images/estabilidad_lineal_weak_mu0_uSi_bajo_umbral.png}
        \subcaption{$\mu=0$, con control adaptativo}
    \end{minipage}\hfill
    \begin{minipage}{0.48\textwidth}
        \centering
        \includegraphics[width=\textwidth]{images/estabilidad_lineal_weak_mu1_uSi_bajo_umbral.png}
        \subcaption{$\mu=1$, con control adaptativo}
    \end{minipage}
    \caption{Espectros lineales para Allee débil con control adaptativo. El control mantiene la inestabilidad mientras concentra la dinámica alrededor de estados cercanos a la extinción ($c\approx 0$, $i\approx 1$).}
    \label{fig:stability_weak_control}
\end{figure}

\section{Allee fuerte}

El sistema con efecto fuerte y difusión se linealiza análogamente:
\begin{align}
    \frac{dc}{dt} &= D_c \nabla^2 c + r_c c \left(\frac{c}{a} - 1\right)\left(\frac{c - a}{k_c - a}\right) - \alpha c s^2 - \beta c i^2 - \mu ( \gamma s^2 + \eta i^2), \\
    \frac{ds}{dt} &= D_s \nabla^2 s + r_s s \left(1 - \frac{s}{k_s} \right) - \gamma s c^2 + \delta s i^2 - \frac{\mu}{2} \alpha s c^2, \\
    \frac{di}{dt} &= D_i \nabla^2 i + r_i i \left(1 - \frac{i}{k_i} \right) - \eta c^2 i + \delta s^2 i - \frac{\mu}{2} \beta i c^2. 
\end{align}

\subsection{Casos sin control adaptativo}

Se presentan los espectros para los 3 escenarios de Allee fuerte sin control adaptativo:

\begin{figure}[H]
    \centering
    \begin{minipage}{0.48\textwidth}
        \centering
        \includegraphics[width=\textwidth]{images/estabilidad_lineal_strong_mu0_uNo_bajo_umbral.png}
        \subcaption{$\mu=0$, bajo umbral}
    \end{minipage}\hfill
    \begin{minipage}{0.48\textwidth}
        \centering
        \includegraphics[width=\textwidth]{images/estabilidad_lineal_strong_mu1_uNo_bajo_umbral.png}
        \subcaption{$\mu=1$, bajo umbral}
    \end{minipage}\\[0.5em]
    \begin{minipage}{0.48\textwidth}
        \centering
        \includegraphics[width=\textwidth]{images/estabilidad_lineal_strong_mu1_uNo_sobre_umbral.png}
        \subcaption{$\mu=1$, sobre umbral}
    \end{minipage}
    \caption{Espectros lineales para Allee fuerte sin control adaptativo. El efecto Allee fuerte produce barreras más pronunciadas; $\mu=1$ reduce la banda de modos inestables en altas frecuencias sin eliminar la inestabilidad en $k\approx 0$.}
    \label{fig:stability_strong_no_control}
\end{figure}

\subsection{Casos con control adaptativo}

Se presentan los espectros para los 2 escenarios de Allee fuerte con control adaptativo:

\begin{figure}[H]
    \centering
    \begin{minipage}{0.48\textwidth}
        \centering
        \includegraphics[width=\textwidth]{images/estabilidad_lineal_strong_mu0_uSi_bajo_umbral.png}
        \subcaption{$\mu=0$, con control adaptativo}
    \end{minipage}\hfill
    \begin{minipage}{0.48\textwidth}
        \centering
        \includegraphics[width=\textwidth]{images/estabilidad_lineal_strong_mu1_uSi_bajo_umbral.png}
        \subcaption{$\mu=1$, con control adaptativo}
    \end{minipage}
    \caption{Espectros lineales para Allee fuerte con control adaptativo. El control colapsa los equilibrios hacia estados cercanos a la extinción manteniendo la inestabilidad, similar al caso de Allee débil.}
    \label{fig:stability_strong_control}
\end{figure}

\section{Resumen comparativo}

En todos los 11 escenarios analizados:
\begin{itemize}
    \item Todos los estados estacionarios son inestables (Re $\lambda_{\max}>0$).
    \item El régimen $\mu=1$ atenúa modos de alta frecuencia pero no estabiliza el sistema en $k\approx 0$.
    \item El control adaptativo mantiene la inestabilidad mientras concentra la dinámica alrededor de estados cercanos a la extinción ($c\approx 0$, $i\approx 1$).
    \item Las partes imaginarias son nulas o constantes en la mayoría de los casos, indicando crecimiento exponencial sin oscilaciones.
    \item El efecto Allee fuerte produce barreras más pronunciadas comparado con el efecto débil.
\end{itemize}



















